\ansclass{15}{Performance Engineering}
\ans{1} (a) $N = X R = 1\,800 \times 0.12 =$ \final{216 holds in flight}. (b) $U = X S = 2\,000 \times 0.0004 =$ \final{80\% busy}. (c) $0.99 \times 20 + 0.01 \times 10\,000 = 19.8 + 100 \approx$ \final{120 ms}: the same mean hides one fan in a hundred waiting 10 s. Report \final{percentiles from histograms} (p50, p95, p99) per endpoint. (d) $2\,000 \times 0.15 =$ \final{300}.
\careful{reading a mean as ``what a typical fan sees''; and averaging percentiles of several pods or runs.}
\ans{2} (a) \final{Smoke}: 1--2 VUs, 1--2 min, zero errors. (b) \final{Load}: expected peak, 15--60 min at the plateau. (c) \final{Soak}: about the daily average for 4--24 h; trends must stay flat. (d) \final{Spike}: baseline to 1--2$\times$ peak in seconds, 5--15 min; tests autoscaling and the waiting room. (e) \final{Breakpoint}: ramp far above the peak until the SLO breaks (abort); the output is a capacity number. (f) \final{Stress}: 1.5$\times$ and 2$\times$ the peak, then back down.
\careful{running only a 10-minute load test: it never reaches the slow failures (pool exhaustion, leaks) nor the 10:00 ramp.}
\ans{3} (a) $X = 200/(R + 1)$: \final{182, 133 and 67 req/s}. (b) $(182 - 67)/182 \approx$ \final{63\% less load} exactly when the system was struggling. An open test keeps sending \final{182 req/s}: at $R = 2$ s the system cannot keep up and \final{the queue grows without bound}, as with real fans. (c) A \final{fixed population}: a partner batch job with 20 workers, a call centre. (d) $1\,000 \times 3.5 =$ \final{3\,500 VUs}. \texttt{dropped\_iterations} means the tester had \final{too few VUs} to start iterations on schedule: the offered load was lower than planned and the run is invalid.
\careful{simulating internet traffic with a fixed number of VUs: the test politely slows down with the system and never builds the queue fans build.}
\ans{4} (a) $60 \times 100 =$ \final{6\,000 samples} in the open model; the closed loop loses $\approx 1.5 \times 100 = 150$ per freeze, \final{$\approx$ 600} in total (the class's run recorded 5\,405). (b) $4 \times 150 / 6\,000 \approx$ \final{10\%} of the samples are slow (the class's p90 is $\approx$ 150 ms). The 60 slowest are the 15 first requests of each freeze, which waited from 1.50 s down to 1.36 s: \final{p99 $\approx$ 1.4 s} (the class computes 1\,366 ms). (c) $4/5\,400 \approx 0.07\%$ of slow samples: \final{p99 $=$ 1 ms}. The closed loop \final{passes} the threshold; the open measurement \final{fails} it, and \final{fans experience the open one}. (d) \final{Open-model generators} (k6 \texttt{constant-arrival-rate}, wrk2) and \final{measuring from the intended start time}; count timeouts at their full duration.
\careful{trusting a perfect p99 from a tool that waits for each response: it recorded the stall once instead of 150 times.}
\ans{5} (a) \final{Knee $\approx$ 1\,100 req/s} (p99 237 ms, the last step within the SLO); \final{maximum $\approx$ 1\,230--1\,250 req/s}. (b) Past saturation, requests time out (5 s) and the work spent on them is wasted, retries add load, contention grows: \final{retrograde} throughput. (c) \final{The knee}: $1\,100/4 =$ \final{275 holds/s per pod}. (d) $180/1\,500 \approx$ \final{120 ms} of pool wait per request. (e) p99 high with CPU low $\to$ requests \final{waiting}: pool 10/10 with a queue, sessions in \texttt{Lock: tuple} $\to$ \final{row-lock contention}, the lock and the connection held during the 40 ms pricing call $\to$ call pricing \final{before \texttt{BEGIN}}, \final{\texttt{FOR UPDATE SKIP LOCKED}}, a waiting room. \final{More pods do not help}: they add waiters for the same rows and connections. (f) $1\,500 \times 0.046 \approx$ \final{69} busy connections inside; $1\,500 \times 0.006 =$ \final{9} before.
\careful{growing the pool to fix pool waits: the database rows are the bottleneck, and more connections only add contention.}
\ans{6} (a) $4\,500/(275 \times 0.75) = 21.8 \to$ \final{22 pods}; $22 \times 3/2 = 33 \to$ \final{33 pods} (a multiple of 3). (b) $3\,000/206.25 = 14.5 \to 15$; $\times 3/2 = 22.5 \to 23 \to$ \final{24 pods}. \texttt{order}: $3\,000/112.5 = 26.7 \to 27 \to 40.5 \to$ \final{42}; with the waiting room $2\,000/112.5 = 17.8 \to 18 \to$ \final{27}. The waiting room cuts these services by about a third. (c) $33 \times 10 =$ \final{330 client connections}, but only $4\,500 \times 0.006 =$ \final{27 active}: put \final{PgBouncer} (transaction pooling) in between. (d) $(4\,500 - 1\,100) \times 120 =$ \final{408\,000 requests}. \final{Pre-scale at 09:45} (with the pools already open) and admit through the waiting room.
\careful{planning with the maximum throughput (1\,250) instead of the knee, or forgetting that 30 more pods mean 300 more database connections.}
\ans{7} (a) A good drawing has the load generator \final{outside} the cluster (not sharing its CPU), the real path (CDN, gateway, services, cache, database with a production-sized snapshot), the payment provider \final{simulated with realistic latency} and RED/USE metrics and traces collected. (b) \final{45\,000 req/s} ($30\,000 \times 1.5$); an \final{open model} ramping 0 $\to$ peak in 30 s; Zipf-like event popularity, 50\,000 synthetic fans, a 60\,000-seat snapshot; payment mock at \final{300 $\pm$ 100 ms}, not 1 ms. (c) $(141 - 120)/120 = 17.5\% > 15\%$: \final{fail}; the largest passing p95 is \final{138 ms}. (d) Thresholds pass in \final{3 consecutive runs}, \texttt{dropped\_iterations} $=$ 0, and \final{$\ge$ 30\% headroom} at the knee.
\careful{testing against tiny data and an instant mock: every query hits a warm cache and the slow dependency never shows up.}
